\documentclass[UTF8]{ctexart}

\usepackage[T1]{fontenc}
\usepackage{lmodern} 
\usepackage{fix-cm}    
\usepackage{microtype} 


\usepackage{listings}
\usepackage{xcolor}
\usepackage{booktabs}
\usepackage{array}
\usepackage{amsmath}
\usepackage{geometry}
\geometry{a4paper, margin=2.5cm}

\lstset{
  language=Python,
  basicstyle=\small\ttfamily,
  keywordstyle=\color{blue},
  commentstyle=\color{gray},
  stringstyle=\color{red},
  showstringspaces=false,
  breaklines=true,
  frame=single,
  numbers=left,
  numberstyle=\tiny\color{gray},
  tabsize=4,
  columns=flexible,
}

\begin{document}

\begin{center}
{\LARGE\textbf{人工智能实验报告}}\\[8pt]
{\Large\textbf{实验一} \quad 知识表示、推理与搜索}
\end{center}

\vspace{8pt}

\begin{center}
\begin{tabular}{
    >{\raggedright\arraybackslash}p{4.8cm}
    >{\centering\arraybackslash}p{4.8cm}
    >{\raggedleft\arraybackslash}p{4.8cm}
}

学院：\underline{\makebox[3.6cm][c]{计算机与通信工程学院}} &
专业：\underline{\makebox[3.6cm][c]{计算机科学与技术}} &
班级：\underline{\makebox[3.6cm][c]{计 234 班}} \\
[6pt]
姓名：\underline{\makebox[3.6cm][c]{秦瑞}} &
学号：\underline{\makebox[3.6cm][c]{U202342492}} &
日期：\underline{\makebox[3.6cm][c]{2026 年 6 月 19 日}} \\
\end{tabular}
\end{center}

\vspace{12pt}

\section{实验目标}

\begin{enumerate}
  \item 熟练掌握知识表示的多种基本方法，包括状态空间法、产生式系统等，并能根据不同问题灵活选择和运用。
  \item 准确实现经典搜索算法，如广度优先搜索（BFS）、深度优先搜索（DFS）、A* 算法（启发式函数）等，理解其算法原理和执行过程。
  \item 深入分析不同搜索策略在效率方面的差异，通过实验结果进行可视化、对比和评估。
\end{enumerate}

\section{实验内容}

\subsection{任务分析}

本次实验包含两个经典人工智能搜索问题：

\textbf{（1）八数码问题（Eight Puzzle）}

给定一个 $3 \times 3$ 的棋盘，其中包含数字 $1$ 至 $8$ 和一个空格（用 $0$ 表示）。通过移动空格相邻的数字块，将初始状态变换为目标状态。初始状态与目标状态如下：

\[
\text{初始状态：}\begin{pmatrix} 1 & 2 & 3 \\ 4 & 0 & 6 \\ 7 & 5 & 8 \end{pmatrix}
\quad\Longrightarrow\quad
\text{目标状态：}\begin{pmatrix} 1 & 2 & 3 \\ 4 & 5 & 6 \\ 7 & 8 & 0 \end{pmatrix}
\]

状态可用九元组表示，如 $(1,2,3,4,0,6,7,5,8)$。每次操作为空格与上下左右相邻位置交换，产生新状态。

\textbf{（2）传教士与食人魔问题（Missionaries and Cannibals）}

河的一侧有 $3$ 名传教士和 $3$ 名食人魔，需要乘船渡河到对岸。约束条件为：
\begin{itemize}
  \item 船每次最多载 $2$ 人，最少载 $1$ 人；
  \item 任意一侧河岸上，若传教士人数大于 $0$，则传教士人数不能少于食人魔人数。
\end{itemize}
状态用三元组 $(M, C, B)$ 表示，其中 $M$ 为左岸传教士人数，$C$ 为左岸食人魔人数，$B$ 为船的位置（$1$ 在左岸，$0$ 在右岸）。初始状态 $(3,3,1)$，目标状态 $(0,0,0)$。

\subsection{实现工具}

\begin{itemize}
  \item \textbf{编程语言}：Python 3
  \item \textbf{核心数据结构}：\texttt{deque}（队列，用于 BFS）、\texttt{list} 模拟栈（用于 DFS）、\texttt{heapq} 优先队列（用于 A*）
  \item \textbf{时间度量}：\texttt{time.perf\_counter()}，提供高精度计时
  \item \textbf{状态表示}：Python 元组（\texttt{tuple}），支持哈希，便于存入 \texttt{set} 和 \texttt{dict}
\end{itemize}

\subsection{实现方案}

整体方案分为三层：

\begin{enumerate}
  \item \textbf{搜索算法层}（\texttt{Search.py}）：实现 BFS、DFS、A* 三种通用搜索算法，接收起始状态、目标状态、邻居函数作为参数，返回搜索结果（是否找到、路径、步数、扩展节点数、耗时等）。
  \item \textbf{问题建模层}（\texttt{EightPuzzle.py}、\texttt{MissionaryCannibal.py}）：为每个问题定义状态表示、邻居生成函数和启发式函数。
  \item \textbf{主程序}：组合调用上述模块，执行搜索并输出结果。
\end{enumerate}

搜索算法均采用统一的 \texttt{SearchResult} 结构返回结果，包含：\texttt{found}（是否找到）、\texttt{path}（路径）、\texttt{steps}（步数）、\texttt{expanded}（扩展节点数）、\texttt{time\_s}（耗时）。

\subsection{实现内容}

\subsubsection{通用搜索算法}

\textbf{广度优先搜索（BFS）}从起始状态开始，逐层扩展节点，使用队列管理待扩展节点。保证找到的路径为最短路径。

\begin{lstlisting}[caption={BFS 算法核心逻辑}]
def BreadthFirstSearch(start, goal, neighbors_fn):
    queue = deque([start])
    came_from = {start: None}
    visited = {start}
    while queue:
        state = queue.popleft()
        if state == goal:
            return SearchResult(True, ReconstructPath(came_from, state), ...)
        for nb in neighbors_fn(state):
            if nb not in visited:
                visited.add(nb)
                came_from[nb] = state
                queue.append(nb)
    return SearchResult(False, None, ...)
\end{lstlisting}

\textbf{深度优先搜索（DFS）}使用栈结构，优先沿一条路径深入搜索到底，再回溯。实现中使用带深度信息的栈元素，并记录每个状态的访问深度，允许在更浅深度再次访问同一状态。

\begin{lstlisting}[caption={DFS 算法核心逻辑}]
def DepthFirstSearch(start, goal, neighbors_fn):
    stack = [(start, 0, None)]
    came_from = {start: None}
    visited_depth = {start: 0}
    while stack:
        state, depth, parent = stack.pop()
        if state == goal:
            return SearchResult(True, ReconstructPath(came_from, state), ...)
        for nb in reversed(neighbors_fn(state)):
            nd = depth + 1
            if nb in visited_depth and visited_depth[nb] <= nd:
                continue
            visited_depth[nb] = nd
            came_from[nb] = state
            stack.append((nb, nd, state))
\end{lstlisting}

\textbf{A* 搜索算法}结合实际代价 $g(n)$ 和启发式估计 $h(n)$，以 $f(n) = g(n) + h(n)$ 为优先级进行搜索。使用最小堆维护开放列表。

\begin{lstlisting}[caption={A* 算法核心逻辑}]
def AStarSearch(start, goal, neighbors_fn, heuristic_fn):
    open_heap = []
    g_score = {start: 0}
    f_score = {start: heuristic_fn(start)}
    heapq.heappush(open_heap, (f_score[start], g_score[start], 0, start))
    came_from = {start: None}
    closed = set()
    while open_heap:
        _, g, _, state = heapq.heappop(open_heap)
        if state in closed:
            continue
        if state == goal:
            return SearchResult(True, ReconstructPath(came_from, state), ...)
        closed.add(state)
        for nb in neighbors_fn(state):
            tentative_g = g + 1
            if tentative_g < g_score.get(nb, float("inf")):
                came_from[nb] = state
                g_score[nb] = tentative_g
                heapq.heappush(open_heap,
                    (tentative_g + heuristic_fn(nb), tentative_g, counter, nb))
\end{lstlisting}

\subsubsection{八数码问题建模}

状态用九元组表示，通过空格（$0$）的位置确定可移动方向。启发式函数采用\textbf{曼哈顿距离}，即每个数字块到其目标位置的行列距离之和（忽略空格）：

\[
h(n) = \sum_{i=1}^{8} |r_i - r_i^{*}| + |c_i - c_i^{*}|
\]

\begin{lstlisting}[caption={八数码问题邻居生成与启发式函数}]
def GetNeighbors(state):
    zero_idx = state.index(0)
    r, c = IndexToMatrix(zero_idx)
    moves = []
    for dr, dc in ((-1,0),(1,0),(0,-1),(0,1)):
        nr, nc = r + dr, c + dc
        if 0 <= nr < 3 and 0 <= nc < 3:
            lst = list(state)
            lst[zero_idx], lst[MatrixToIndex(nr,nc)] = lst[MatrixToIndex(nr,nc)], lst[zero_idx]
            moves.append(tuple(lst))
    return moves

def ManhattanDistance(state, goal=GOAL_STATE):
    pos_goal = {val: i for i, val in enumerate(goal)}
    total = 0
    for i, val in enumerate(state):
        if val == 0: continue
        gi = pos_goal[val]
        r1,c1 = IndexToMatrix(i)
        r2,c2 = IndexToMatrix(gi)
        total += abs(r1-r2) + abs(c1-c2)
    return total
\end{lstlisting}

\subsubsection{传教士与食人魔问题建模}

状态用三元组 $(M, C, B)$ 表示，每次渡河选择 $1$ 或 $2$ 人的组合（共 $5$ 种合法移动）。合法性检验确保两岸传教士人数均不少于食人魔人数（或传教士为零时不受约束）。

启发式函数设计为剩余过河人数的一半（向上取整），因为每次最多载 $2$ 人：

\[
h(n) = \left\lceil \frac{M + C}{2} \right\rceil
\]

\begin{lstlisting}[caption={传教士与食人魔问题启发式函数}]
def Heuristic(state):
    missionary, cannibal, boat = state
    remaining = missionary + cannibal
    return (remaining + 1) // 2
\end{lstlisting}

\subsubsection{实验结果与分析}

\textbf{（1）八数码问题搜索结果}

对初始状态 $(1,2,3,4,0,6,7,5,8)$ 进行搜索，结果如下表所示：

\begin{table}[htbp]
\centering
\caption{八数码问题三种搜索算法对比}
\begin{tabular}{lcccc}
\toprule
算法 & 是否找到 & 步数 & 扩展节点数 & 耗时 (s) \\
\midrule
BFS & 是 & 2 & 9 & 0.0 \\
DFS & 是 & 2 & 14 & 0.0001 \\
A* & 是 & 2 & 3 & 0.0001 \\
\bottomrule
\end{tabular}
\end{table}

A* 搜索的解路径为：

\[
(1,2,3,4,\mathbf{0},6,7,5,8) \rightarrow (1,2,3,4,5,6,7,\mathbf{0},8) \rightarrow (1,2,3,4,5,6,7,8,\mathbf{0})
\]

\textbf{分析}：
\begin{itemize}
  \item 三种算法均在 $2$ 步内找到解，但扩展节点数差异显著。A* 仅扩展 $3$ 个节点，BFS 扩展 $9$ 个，DFS 扩展 $14$ 个。
  \item A* 通过曼哈顿距离启发式有效引导搜索方向，大幅减少了不必要的节点扩展。
  \item DFS 虽然时间开销与 A* 相近（均约为 $0.0001$s），但扩展了更多节点，且最大搜索前沿达到 $5$，内存占用更高。
  \item BFS 扩展节点数居中（$9$ 个），但保证了最短路径。
\end{itemize}

\textbf{（2）传教士与食人魔问题搜索结果}

对初始状态 $(3,3,1)$ 进行搜索，结果如下表所示：

\begin{table}[htbp]
\centering
\caption{传教士与食人魔问题搜索算法对比}
\begin{tabular}{lcccc}
\toprule
算法 & 是否找到 & 步数 & 扩展节点数 & 耗时 (s) \\
\midrule
BFS & 是 & 11 & 15 & 0.0 \\
A* & 是 & 11 & 15 & 0.0001 \\
\bottomrule
\end{tabular}
\end{table}

A* 搜索的完整解路径如下：

\begin{table}[htbp]
\centering
\caption{传教士与食人魔问题 A* 解路径}
\begin{tabular}{ccccccccccccc}
\toprule
步骤 & 1 & 2 & 3 & 4 & 5 & 6 & 7 & 8 & 9 & 10 & 11 & 12 \\
\midrule
$M$ & 3 & 3 & 3 & 3 & 3 & 1 & 2 & 0 & 0 & 0 & 1 & 0 \\
$C$ & 3 & 1 & 2 & 0 & 1 & 1 & 2 & 2 & 3 & 1 & 1 & 0 \\
$B$ & 1 & 0 & 1 & 0 & 1 & 0 & 1 & 0 & 1 & 0 & 1 & 0 \\
\bottomrule
\end{tabular}
\end{table}

\textbf{分析}：
\begin{itemize}
  \item BFS 和 A* 均找到步数为 $11$ 的解，且扩展节点数相同（$15$ 个）。
  \item 对于此问题，由于状态空间较小（合法状态仅约 $20$ 个），启发式函数的引导效果有限，BFS 与 A* 性能接近。
  \item 启发式函数 $h(n) = \lceil(M+C)/2\rceil$ 给出了每次渡河最多 $2$ 人的下界估计，保证了 A* 的最优性。
  \item 该问题的解路径较长（$11$ 步），反映了传教士与食人魔问题固有的约束复杂性——需要反复来回运输以维持安全约束。
\end{itemize}

\textbf{（3）三种搜索算法综合对比}

\begin{table}[htbp]
\centering
\caption{算法特性对比总结}
\begin{tabular}{lccc}
\toprule
特性 & BFS & DFS & A* \\
\midrule
最优性 & 保证最短路径 & 不保证 & 保证（启发函数一致时） \\
完备性 & 保证 & 不保证（可能陷入无限深度） & 保证 \\
空间复杂度 & $O(b^d)$ & $O(b \cdot m)$ & $O(b^d)$ \\
时间复杂度 & $O(b^d)$ & $O(b^m)$ & $O(b^d)$（依赖启发函数质量） \\
\bottomrule
\end{tabular}
\end{table}

其中 $b$ 为分支因子，$d$ 为解的深度，$m$ 为搜索最大深度。

\section{实验总结}

\textbf{结果分析}：实验结果表明，在八数码问题中，A* 算法凭借曼哈顿距离启发式，仅需扩展 $3$ 个节点即可找到最优解，相比 BFS（$9$ 个节点）和 DFS（$14$ 个节点）具有显著的效率优势。在传教士与食人魔问题中，由于状态空间较小，BFS 与 A* 表现相近。

\textbf{算法性能对比}：BFS 保证最短路径但空间开销大；DFS 空间效率高但不保证最优解，且可能陷入深层搜索；A* 在保证最优性的前提下，通过启发式函数有效缩减搜索空间，是实际应用中最常用的搜索策略。

\textbf{启发式函数的作用}：启发式函数为 A* 提供了对目标距离的估计，使搜索更具方向性。曼哈顿距离对八数码问题效果显著，因为它是可采纳的（不会高估实际代价）且一致的。传教士与食人魔问题的启发式虽较简单，但也有效引导了搜索。

\textbf{遇到的问题与解决}：DFS 实现中需要处理状态的重复访问问题。通过记录每个状态的访问深度，允许在更浅深度再次访问同一状态，避免遗漏更优路径。

\textbf{思考与启迪}：本实验加深了对经典搜索算法原理的理解。在实际问题中，选择合适的搜索策略和启发式函数至关重要——A* 算法的成功很大程度上取决于启发式函数的质量。未来可进一步研究迭代加深 A*（IDA*）等内存优化算法，以及将搜索技术应用于更复杂的问题域。

\end{document}
